\documentclass[11pt,a4paper]{article}
\usepackage[margin=1in]{geometry}
\usepackage[T1]{fontenc}
\usepackage{graphicx}
\usepackage{booktabs}
\usepackage{float}
\usepackage{hyperref}
\usepackage{caption}
\usepackage{enumitem}
\usepackage{xurl}
\usepackage{array}
\emergencystretch=3em

% Figure if the file exists, otherwise a visible placeholder box.
\newcommand{\safefig}[3]{%
  \begin{figure}[H]\centering
  \IfFileExists{../reports/figures/#1}{\includegraphics[width=#2\linewidth]{../reports/figures/#1}}{%
    \fbox{\parbox{0.8\linewidth}{\centering\vspace{1.2cm}\detokenize{#1} not found: run the EDA scripts first.\vspace{1.2cm}}}}%
  \caption{#3}\end{figure}}

% Run LaTeX from paper/. Pilot PNGs belong in figures/pilot_eda/.
\newcommand{\pilotfig}[3]{%
  \begin{figure}[H]\centering
  \IfFileExists{figures/pilot_eda/#1}{%
    \includegraphics[width=#2\linewidth]{figures/pilot_eda/#1}%
  }{%
    \fbox{\parbox{0.8\linewidth}{\centering\vspace{0.8cm}%
    Missing pilot figure: \texttt{\detokenize{#1}}\vspace{0.8cm}}}%
  }%
  \caption{#3}\end{figure}}

\newcommand{\modelingfig}[3]{%
  \begin{figure}[H]\centering
  \includegraphics[width=#2\linewidth]{figures/modeling/#1}%
  \caption{#3}
  \end{figure}}

\title{Federated Learning for EEG-Based Working-Memory Decoding\\
\large Model Architecture and Empirical Evaluation (Deliverable 2)}
\author{Eman Adil \and Dua Sarfraz \and Ayza Ahmad\\[3pt] {Artificial Intelligence}}
\date{\today}

\begin{document}
\maketitle

\begin{abstract}
We investigate EEG-based working-memory load decoding with centralized,
federated and personalized learning on OpenNeuro \texttt{ds003655}.
A two-participant preprocessing pilot was expanded to 24 participants.
The primary dataset comprised 2,878 trials and 57 delay-period band-power
features per trial. Three neural-network training strategies were
evaluated on 359 matched trials from four held-out participants.
Centralized, FedAvg and personalized FedAvg macro-F1 scores were 0.3036,
0.3276 and 0.3397, respectively. All three ROC-AUC scores were near 0.50,
indicating limited discrimination on this split.
\end{abstract}

\tableofcontents
\newpage

%===========================================================
\section{Phase 1: Problem Definition}

\subsection{Problem Statement}
Working memory is the short-term maintenance and manipulation of information, and it is central to reasoning and
learning. EEG measures brain activity non-invasively and cheaply, so decoding working-memory demand from EEG is of
interest for cognitive-state monitoring and brain--computer interfaces. Two practical obstacles limit this:
\begin{enumerate}[leftmargin=*]
  \item \textbf{Inter-subject variability.} EEG differs strongly between people (anatomy, electrode contact, noise,
  individual rhythm frequencies), so a model trained on some people often degrades on others, and data from many
  participants is needed.
  \item \textbf{Privacy and data silos.} EEG is sensitive personal data. Pooling recordings from different labs or
  clinics is often restricted by consent and regulation.
\end{enumerate}
Federated learning trains a shared model by exchanging model updates instead of raw recordings, which addresses the
second obstacle. However, each participant's data is distributed differently (non-IID), which is exactly the first
obstacle and is known to make federated optimisation harder \cite{mcmahan2017}. This project studies that
interaction for EEG working-memory decoding.

\subsection{Research Questions}
\begin{enumerate}[leftmargin=*]
  \item How does a federated model compare with a centralized baseline for working-memory condition decoding?
  \item How sensitive are results to participant-level distribution differences?
  \item Does a clearly defined personalization strategy improve cross-subject generalization?
\end{enumerate}
These are planned experiments, not results at this stage.

\subsection{Target Audience and Research Gap}
This is an original research (R\&D) project. The intended audience is researchers in neuro-engineering,
brain--computer interfaces and privacy-preserving machine learning. The gap is a reproducible, small-scale comparison
of centralized, federated and personalized learning for EEG working-memory decoding under non-IID, cross-subject
conditions, in which data quality and label validity are examined rather than assumed.

\textbf{Scope.} The initial pilot uses two participants. The later modeling experiment expands to 24 participants.

\subsection{Justification for Machine Learning}
A rule-based program would need hand-written rules from EEG measurements to working-memory demand. This is unsuitable
because:
\begin{itemize}[leftmargin=*]
  \item \textbf{High-dimensional, weak, distributed signal.} Each trial is a 19-channel time series at 500\,Hz;
  memory-load effects appear as subtle spectral changes spread over channels, with no single threshold separating
  conditions.
  \item \textbf{Subject-specific patterns.} The same state looks different across people, so fixed rules do not transfer,
  whereas learned models can adapt to these differences.
  \item \textbf{Low signal-to-noise ratio.} Artifacts (eye movement, muscle activity, amplifier saturation; see
  Section~\ref{sec:quality}) require statistical models that weigh evidence over many features.
  \item \textbf{The research question itself is about learning.} Federated optimisation and personalization are
  properties of learning algorithms and have no rule-based counterpart.
\end{itemize}

%===========================================================
\section{Phase 2: Data Acquisition}

\subsection{Dataset and Collection}
The data are from OpenNeuro dataset \texttt{ds003655}, task \texttt{VerbalWorkingMemory}
(\url{https://openneuro.org/datasets/ds003655}), a modified Sternberg task with retention (R) and manipulation (M)
conditions at set sizes 5, 6 and 7, giving the conditions 5R, 6R, 7R, 5M, 6M and 7M. Two participants,
\texttt{sub-001} and \texttt{sub-002}, were downloaded from the public repository. The local inventory
(\texttt{reports/download\_inventory.csv}, produced by \texttt{01\_check\_downloads.py}) records eight files per
participant and about 269.05\,MB in total. Raw files are excluded from Git through \texttt{.gitignore}. No download
API log was captured.

\begin{table}[H]\centering
\caption{Recording metadata.}
\begin{tabular}{lcc}
\toprule
Property & sub-001 & sub-002 \\
\midrule
Sampling rate & 500 Hz & 500 Hz \\
Channels recorded & 21 & 21 \\
EEG channels & 19 & 19 \\
EOG channels & EOGv, EOGh & EOGv, EOGh \\
Duration & 3077 s & 3035 s \\
\bottomrule
\end{tabular}
\end{table}
The EEG channels are Fp1, Fp2, F7, F3, Fz, F4, F8, T3, C3, Cz, C4, T4, T5, P3, Pz, P4, T6, O1, O2.

\subsection{Trial Structure}
\texttt{04\_build\_trial\_table.py} builds \texttt{reports/trial\_table.csv} from the event markers: 240 trials
(120 per participant), 20 per condition per participant. Each row holds the onsets of the baseline, encoding, delay
and probe phases, the response information and quality flags. The delay-to-probe interval is about 6.7\,s, but all
scripts use the actual event onsets.

%===========================================================
\section{Phase 2: Validation and Data-Centric Cleaning}
\label{sec:quality}

\subsection{File and Event Validation}
Both recordings loaded correctly, with 120 trial starts and 20 baseline, encoding, delay and probe markers per
condition. The event tables had no exact duplicate rows and no onsets outside the recording. These checks confirm
structure, not scientific validity.

\subsection{Response-Label Uncertainty}
There are 244 response markers for 240 probes; some are late, duplicated or ambiguous. The trial table flags eight trials
for response ambiguity (6 in sub-001, 2 in sub-002), and the later probe-based audit marks 225 of 240 rows as PASS and 15
as REVIEW; the two counts come from different checks and are to be reconciled. Late responses are judged with a five-second heuristic that is not a
validated rule. Uncertain cases are kept, not silently resolved, and the event-to-label mapping is provisional; no
supervised model will be trained before the labels are validated.

\subsection{Channel-Type Correction}
MNE initially labelled \texttt{EOGv} and \texttt{EOGh} as EEG. They are now set to type \texttt{eog} before channel selection,
so features use only the 19 EEG channels (\texttt{12\_inspect\_channel\_types.py}, the Colab pilot notebook).

\subsection{Initial Signal-Quality Findings}
\textbf{sub-001 boundary values.} In all 21 channels the minimum value reaches about $-10{,}239.69\,\mu$V (maximum
absolute value 10{,}239.69\,$\mu$V). Between 332.5 and 333.5\,s each channel has many samples near this lower boundary,
including consecutive runs of about 0.48--0.52\,s, and the same boundary was not observed in sub-002 in that interval.
This is suspicious for clipping, saturation or a data-representation issue, but \textbf{the cause is not confirmed}
(\texttt{reports/signal\_quality\_decisions.md}).

\textbf{Amplitude screening before correction.} Each trial window was flagged if its maximum channel peak-to-peak or absolute amplitude
exceeds a heuristic threshold of 500\,$\mu$V.
\begin{table}[H]\centering
\caption{Heuristic amplitude flags per trial window (not validated rejection rules).}
\begin{tabular}{lccc}
\toprule
Participant & No flag & Peak-to-peak only & Both flags \\
\midrule
sub-001 & 112 & 6 & 2 \\
sub-002 & 0 & 0 & 120 \\
\midrule
Total & 112 & 6 & 122 \\
\bottomrule
\end{tabular}
\end{table}
\begin{table}[H]\centering
\caption{Window amplitude by flag group ($\mu$V).}
\begin{tabular}{lccc}
\toprule
Group & Median p--p & Mean p--p & Median abs.\ amp. \\
\midrule
No flag & 406.66 & 411.35 & 258.31 \\
Peak-to-peak only & 524.28 & 550.06 & 371.22 \\
Both flags & 2550.53 & 2588.39 & 1325.66 \\
\bottomrule
\end{tabular}
\end{table}
Every sub-002 trial is flagged, which does not by itself mean every trial should be discarded. Whole-record
channel ranges in sub-002 are large but vary by channel (e.g.\ Pz 5461.00, T3 5318.31, O1 5506.38\,$\mu$V), and the EOG
channels also have large extrema. To tell a single bad channel from a participant-wide difference, the per-channel,
per-trial audit (\texttt{24\_inspect\_trial\_channel\_amplitudes.py}, output
\texttt{reports/trial\_channel\_amplitude\_audit.csv}) was computed for every delay-to-probe window
(Table~\ref{tab:chanaudit}).
\begin{table}[H]\centering
\caption{Per-trial channel amplitude audit by participant (120 windows each; $\mu$V).}\label{tab:chanaudit}
\begin{tabular}{>{\raggedright\arraybackslash}p{5.4cm}>{\raggedright\arraybackslash}p{4.1cm}>{\raggedright\arraybackslash}p{4.1cm}}
\toprule
Quantity & sub-001 & sub-002 \\
\midrule
Max channel peak-to-peak, median (range) & 410.8 (317.4--1612.2) & 2551.4 (2444.9--2989.5) \\
Max absolute amplitude, median & 260.5 & 1325.7 \\
Median-over-channels peak-to-peak, mean & 256.5 & 613.0 \\
Channels above 500\,$\mu$V p--p per window, mean (range) & 0.24 (0--19) & 13.2 (12--19) \\
Channel with the largest p--p & Fp2 50, T3 45, O2 24, F8 1 trials & T3 in all 120 trials \\
\bottomrule
\end{tabular}
\end{table}
Two observations follow. First, the sub-002 difference is not confined to one channel: the typical channel is itself
about 2.4 times larger than in sub-001, and 12 or more of the 19 EEG channels exceed 500\,$\mu$V in every window. Second,
T3 has the largest range in every sub-002 window, so a channel-specific problem may be superimposed on the general
difference. In sub-001 the flags are rare and concentrated in a few windows (the maximum of 19 channels over threshold
belongs to a single window). The cause of the sub-002 pattern (recording or reference differences, scaling, drift, or
genuine activity) is \textbf{still not established}. Note also that even the unflagged sub-001 windows have a median
channel peak-to-peak range of roughly 250\,$\mu$V, which is large for EEG; we hypothesise that this reflects the
absence of high-pass filtering, but this has not been tested.

\subsection{Signal Correction and Final Pilot Audit}
The pilot notebook cropped each recording from 10\,s before the first delay
to 10\,s after the last probe, within recording bounds. The suspicious
sub-001 interval preceded this crop. Processing applied a 50\,Hz IIR notch,
a 1--40\,Hz fourth-order Butterworth IIR band-pass, and resampling from
500 to 100\,Hz. The linked-earlobe reference was retained.

FastICA was fitted separately to the EEG channels of each participant using
a 0.99 variance fraction, seed 42 and a maximum of 1{,}000 iterations.
Component review used EOG association, scalp maps and time-course inspection.
Component 0 was removed from each participant's own fit. For sub-002,
component 1 was retained after review. Across the 120 delay windows per
participant, the median absolute Fp2 correlation with vertical EOG went from
0.896 to 0.059 for sub-001 and from 0.307 to 0.081 for sub-002.

\begin{table}[H]\centering
\caption{Final peak-to-peak screening of corrected delay windows. The amplitude
columns summarize the maximum across EEG channels in each trial.}
\begin{tabular}{lrrrr}
\toprule
Participant & Trials & Median ($\mu$V) & Highest ($\mu$V) & Trials $>500\,\mu$V \\
\midrule
sub-001 & 120 & 67.458 & 142.577 & 0 \\
sub-002 & 120 & 66.930 & 249.077 & 0 \\
\bottomrule
\end{tabular}
\end{table}

No channel within a trial fell below the flatness threshold of
0.1\,$\mu$V peak-to-peak. All 240 trials were retained. The results
are recorded in \texttt{final\_pilot\_signal\_quality.csv}.

\subsection{Tabular Cleaning of the Feature Table}
The corrected feature audit found zero missing, infinite or negative band-power
values, zero duplicate participant/trial keys and zero duplicate full rows.
No imputation was needed. For each participant and feature, values outside
the interval from $Q_1-1.5\,\mathrm{IQR}$ to $Q_3+1.5\,\mathrm{IQR}$ were flagged
for review. At least one of 57 features was flagged in 56 trials for sub-001
and 57 for sub-002. All flagged rows were retained. The 15 response-review
flags were kept separately from the feature flags.

\section{Phase 2: Feature Strategy}

\subsection{Spectral Features}
We used the corrected delay-to-probe windows for all 19 EEG channels.
Welch power spectral density \cite{welch1967} was computed with a Hann window,
2\,s segments, 1\,s overlap, constant detrending and density scaling. Signals
were expressed in microvolts before estimating the spectrum. Trapezoidal
integration over theta (4--8\,Hz), alpha (8--13\,Hz) and beta (13--30\,Hz)
included both frequency endpoints. Each trial produced $19\times3=57$
band-power features in $\mu\mathrm{V}^2$. The table
\texttt{eeg\_features\_cleaned\_signal.csv} also records trial timing, condition
and signal-quality metadata.

\subsection{Scaling and Encoding}
For the scaling demonstration, the 57 band-power columns were standardized
with means and population standard deviations estimated from sub-001 and
applied to both participants. There were no constant reference features.
The median feature mean and standard deviation were 0 and 1 for sub-001
and $-0.2907$ and 0.8794 for sub-002. The resulting table is
\texttt{eeg\_features\_scaled\_demo.csv}. Future evaluations will fit scaling
parameters within each training split.

Conditions 5R, 6R, 7R, 5M, 6M and 7M map to nominal IDs 0 through 5,
respectively. Set size (5, 6 or 7) and task type (R or M, encoded as 0 or 1)
are also stored in \texttt{eeg\_features\_encoded.csv}. The input feature
list contains only band powers. Trial identifiers, labels, response audit
fields and quality flags are stored as separate columns.

\section{Exploratory Data Analysis}
\subsection{Exploratory analysis of the two-participant pilot}
Each participant contributed 20 trials to each of the six conditions,
giving 40 trials per condition and 240 in total.

\pilotfig{01_trial_balance.png}{0.80}
{Trial counts by condition and participant.}

For each trial and frequency band, we summarized delay-period power
using the median across 19 EEG channels. Sub-001 had higher median theta
and beta power. Sub-002 had slightly higher median alpha power and a
wider alpha interquartile range.

\pilotfig{02_band_power_distributions.png}{0.98}
{Distribution of trial-level median EEG channel power by band and participant.}

The Pearson heatmap showed groups of positive correlations within
bands. Beta also contained groups of negative correlations between channels.
Correlations between bands varied in sign. The analysis pools trials from
both participants, so recording differences can contribute to these patterns.

\pilotfig{03_feature_correlations.png}{0.74}
{Pearson correlations among 57 unscaled band-power features across 240 pilot trials.}

\subsection{Response-marker review}
\safefig{response_audit_status.png}{0.70}
{Response-marker review status by participant.}
Of 240 probes, 225 passed the structural audit and 15 required review
(eight for sub-001 and seven for sub-002).

\section{Limitations of the Initial Pilot}
\begin{enumerate}[leftmargin=*]
  \item Findings describe two pilot recordings. The expanded cohort is described in Phase 3.
  \item Response accuracy remains provisional for ambiguous markers. The
  condition labels used here come from the experimental task markers.
  \item The cause of the pre-task sub-001 boundary values is unconfirmed.
  Amplitude and flatness thresholds are screening heuristics.
  \item The pilot ICA fits used each participant's full task crop. Future
  evaluations must define the ICA fitting or calibration procedure relative
  to the evaluation split, alongside training-only scaling.
  \item Response-label validation remains open. Phase 3 and Phase 4 describe
  the expanded cohort and the model comparison.
\end{enumerate}

\section{Reproducibility}
The repository contains the acquisition and audit scripts under
\texttt{data\_scripts/}. The Colab notebook
\texttt{notebooks/EEG\_Preprocessing\_Pipeline.ipynb} records the pilot signal
correction, feature extraction, encoding, audit, scaling demonstration and
EDA. Reviewed ICA solutions, corrected recordings and derived tables are
saved to Drive checkpoints. Raw EEG stays outside the repository.
Repository: \url{https://github.com/EmanAdil122/federated-eeg-working-memory}.

\section{Phase 3: Model Architecture and Implementation}

\subsection{Expanded cohort and prediction target}
The two-participant pilot workflow was extended to 24 participants from
OpenNeuro \texttt{ds003655}. The merged table contained 2,880 trials.
Two trials with abrupt residual EEG amplitudes above the 500\,$\mu$V
review threshold were excluded from the primary analysis
(sub-014 trial 62 and sub-020 trial 53), leaving 2,878 trials.
The classifier uses 57 delay-period EEG band-power features and predicts
memory load (5, 6 or 7 letters) across retention (R) and manipulation (M).
The target comes from task-condition markers. Response correctness and
response-review status were excluded from the model inputs.

\subsection{Shared classifier and training strategies}
We implemented a multilayer perceptron in PyTorch with 57 inputs, hidden
layers of 64 and 32 units with ReLU activations, dropout of 0.4 after the
first hidden layer, and three output logits. Training used cross-entropy
and Adam. The features were standardized using the means and standard
deviations estimated from the 16 training participants and applied to
validation and test features. All strategies used the same feature scale.

\modelingfig{eeg_ann_architecture.png}{0.98}
{Shared architecture and the three training strategies.}

\begin{itemize}[leftmargin=*]
  \item \textbf{Centralized ANN.} Pooled the training participants' features
  and trained with mini-batches of 64 for 30 epochs. The checkpoint with
  the best validation macro-F1 was selected.
  \item \textbf{FedAvg ANN.} Simulated 16 participant clients in one runtime.
  Each client started each round from the current global weights, trained
  locally for one epoch, and contributed weights to the server average
  in proportion to its trial count \cite{mcmahan2017}. All clients
  participated in each of the 30 rounds. The architecture and optimizer
  settings were carried over from centralized tuning.
  \item \textbf{Personalized FedAvg ANN.} Initialized from the selected
  global FedAvg checkpoint. For each evaluation participant, the output
  layer was adapted with Adam using 30 labeled calibration trials, the
  first five primary trials in each of the six conditions. The hidden
  layers were frozen. Calibration used full-batch updates, a learning rate
  of 0.001 and 10 epochs selected on validation participants.
\end{itemize}

\subsection{Hyperparameter selection}
Table~\ref{tab:central-tuning} reports the best validation checkpoint for
each centralized configuration. The batch size was held fixed at 64.
Dropout of 0.4 was selected because it gave the highest validation
macro-F1 (0.4010). The FedAvg checkpoint at round 8 reached 0.3014 macro-F1
and 32.29\% validation accuracy on the full validation set.

\begin{table}[H]
\centering
\caption{Centralized ANN tuning. Checkpoints were ranked by validation macro-F1.}
\label{tab:central-tuning}
\small
\begin{tabular}{lrrrrrr}
\toprule
Configuration & Hidden units & Dropout & Learning rate & Epoch & Accuracy & Macro-F1 \\
\midrule
Baseline & 64, 32 & 0.2 & 0.0010 & 10 & 0.3833 & 0.3829 \\
Smaller network & 32, 16 & 0.2 & 0.0010 & 12 & 0.3833 & 0.3791 \\
More dropout & 64, 32 & 0.4 & 0.0010 & 10 & 0.4042 & 0.4010 \\
Lower learning rate & 64, 32 & 0.2 & 0.0005 & 11 & 0.3938 & 0.3796 \\
\bottomrule
\end{tabular}
\end{table}

Personalization was tuned on held-out validation trials after calibration.
Mean participant macro-F1 for 0, 1, 3, 5 and 10 adaptation epochs was
0.2518, 0.2540, 0.2682, 0.2819 and 0.2986, respectively. These scores
use a different validation subset and aggregation from the full-set
FedAvg checkpoint score.

\modelingfig{training_and_validation_curves.png}{0.98}
{Centralized loss and validation accuracy by epoch, and FedAvg client training
loss and validation accuracy by round. Dotted vertical lines mark checkpoints
selected by validation macro-F1.}

The centralized training loss decreased while validation loss increased
in later epochs, consistent with overfitting on this split. FedAvg
validation accuracy stayed near or below the uniform-guess reference
for most rounds.

\section{Phase 4: Training and Empirical Evaluation}

\subsection{Participant split and test protocol}
Participants were assigned to training (16), validation (4) and test (4)
using seed 42. No participant was assigned to more than one split.
Test participants were sub-002, sub-009, sub-014 and sub-022.
Validation results determined model checkpoints and adaptation settings.
Of the test participants' 479 primary trials, 120 were reserved for
personalized calibration (30 per participant). The three models were
scored on the same remaining 359 trials. Centralized and FedAvg weights
were fixed for testing, while personalization used only the labeled
calibration subset for the respective test participant. Calibration
and evaluation trials came from the same recording session.

Accuracy measures overall correctness. Macro precision, recall and F1
give each load class equal weight. Macro one-versus-rest ROC-AUC
summarizes how well predicted probabilities rank each class against the
other two classes, independently of the final class decision.

\subsection{Held-out results}
\begin{table}[H]
\centering
\caption{Final test performance on 359 matched trials from four participants.}
\label{tab:test-performance}
\small
\begin{tabular}{lrrrrr}
\toprule
Model & Accuracy & Macro precision & Macro recall & Macro-F1 & Macro ROC-AUC \\
\midrule
Centralized ANN & 0.3036 & 0.3098 & 0.3035 & 0.3036 & 0.4752 \\
FedAvg ANN & 0.3538 & 0.3573 & 0.3543 & 0.3276 & 0.4750 \\
Personalized FedAvg ANN & 0.3538 & 0.3643 & 0.3541 & 0.3397 & 0.4797 \\
\bottomrule
\end{tabular}
\end{table}

FedAvg achieved higher accuracy and macro-F1 than the centralized model
on this held-out split. Personalization kept accuracy at 35.38\% and
raised macro-F1 from 0.3276 to 0.3397. Accuracy remained close to the
33.33\% uniform-guess reference, and macro ROC-AUC ranged from 0.4750 to
0.4797. These results provide limited evidence for reliable load decoding
or consistent superiority on unseen participants.

\subsection{Error Analysis}
\modelingfig{test_confusion_matrices.png}{0.99}
{Test confusion matrices. Rows are actual loads and columns are predicted loads.
Each cell shows the count and percentage within the actual class.}

The centralized model predicted load 6 on 156 of the 359 trials.
FedAvg predicted load 5 on 183 trials. Load 7 was particularly difficult
for both federated models. FedAvg correctly classified 15 of 120 load-7
trials (12.5\% recall); personalization raised this to 21 of 120 (17.5\%).
Personalization also improved load-6 recall from 37.5\% to 40.0\%, while
load-5 recall dropped from 56.3\% to 48.7\%.

\subsection{Bias Check}
\begin{table}[H]
\centering
\caption{Descriptive test macro-F1 by participant and task.}
\label{tab:bias-check}
\begin{tabular}{lrrr}
\toprule
Group & Centralized & FedAvg & Personalized FedAvg \\
\midrule
sub-002 & 0.350 & 0.340 & 0.371 \\
sub-009 & 0.242 & 0.303 & 0.317 \\
sub-014 & 0.209 & 0.243 & 0.281 \\
sub-022 & 0.241 & 0.270 & 0.290 \\
\midrule
Manipulation (M) & 0.273 & 0.316 & 0.292 \\
Retention (R) & 0.325 & 0.330 & 0.383 \\
\bottomrule
\end{tabular}
\end{table}

Personalization improved participant-level macro-F1 relative to FedAvg
for each of the four test participants. The task-level comparison
differed: retention improved from 0.330 to 0.383, while manipulation
decreased from 0.316 to 0.292. These comparisons describe performance
variation within the available split. Demographic fairness was not
assessed in this experiment.

\subsection{Limitations}
The evaluation uses one split and only four test participants.
ICA correction was fitted separately on each participant's full task
recording before model splitting. This protocol assumes access to the
full recording for preprocessing and does not establish prospective
performance on an incoming trial stream. Personalization uses labeled
trials from each test participant's session, so it measures adaptation
to that session. FedAvg was a single-runtime simulation using a shared
scaler estimated from training participants. It does not demonstrate
a deployed privacy system.

\section{Contribution Table}
\begin{table}[H]
\centering
\caption{Contributions verified in repository history at the time of drafting.}
\small
\begin{tabular}{p{0.20\linewidth}p{0.23\linewidth}p{0.45\linewidth}}
\toprule
Member & Git author & Module and example commit \\
\midrule
Eman Adil & \texttt{EmanAdil122} &
EEG acquisition, validation and feature extraction
(\texttt{6134493}, \texttt{364382c}) \\
Dua Sarfraz & \texttt{duasarfarz3114} &
Feature EDA script and visualization fixes
(\texttt{64e25ca}, \texttt{7dc6a6a}) \\
Ayza Ahmad & \texttt{ayzaahmed27} &
Batch EEG preprocessing and 24-participant quality review
(\texttt{1ccc3f7}, \texttt{53df3a5}) \\
\bottomrule
\end{tabular}
\end{table}
% Before submission, add each member's actual Phase 3/4 module and commit ID.

\begin{thebibliography}{9}
\bibitem{welch1967} P.~Welch, ``The use of fast Fourier transform for the estimation of power spectra,'' \emph{IEEE Trans.\ Audio Electroacoust.}, 15(2):70--73, 1967.
\bibitem{mne2013} A.~Gramfort et al., ``MEG and EEG data analysis with MNE-Python,'' \emph{Front.\ Neurosci.}, 7:267, 2013.
\bibitem{mcmahan2017} B.~McMahan et al., ``Communication-efficient learning of deep networks from decentralized data,'' \emph{AISTATS}, 2017.
\end{thebibliography}
\end{document}